\section{Simulation details}\label{app:simdetails}

This appendix documents the simulation study of \cref{sec:simulation}: software, seeds and runs (\cref{app:sim-run}), configurations (\cref{app:sim-config}), estimator settings (\cref{app:sim-est}), Monte Carlo standard errors (\cref{app:sim-mcse}), exact checks and recomputed illustrations (\cref{app:sim-checks}), further result tables (\cref{app:sim-tables}) and reproduction (\cref{app:sim-repro}). All data are synthetic. Every configuration value below is stored with the results in \texttt{results/sim\_results.json}. The per-replication values are in \texttt{results/raw/E1.json}, \dots, \texttt{results/raw/E5.json}, with the near-proportional design in \texttt{results/raw/E1b.json}.

\subsection{Software, seeds and runs}\label{app:sim-run}

\paragraph{Software}
The experiments are run by \texttt{code/run\_experiments.py}, which calls the package \texttt{code/attsim/}: data generation (\texttt{generate.py}, \texttt{designs.py}), polychoric correlations (\texttt{polychoric.py}, \texttt{bvn.py}), bifactor and correlated-factors fits (\texttt{fit.py}), boundary estimands and intervals (\texttt{boundary.py}), DIF procedures (\texttt{dif.py}), list experiments (\texttt{listexp.py}), MRP (\texttt{mrp.py}) and count indices (\texttt{index.py}). The figures and the tables \texttt{tables/tab-sim-*.tex} are written by \texttt{code/make\_figures.py} from the results file. The recorded environment is Python 3.13.18, NumPy 2.5.3, SciPy 1.18.1 and Matplotlib 3.11.2 on Linux x86\_64 (glibc 2.39). No other packages are used. The unit tests in \texttt{code/tests/} (26 tests, run with \texttt{pytest}) check the generator, the polychoric estimator and its standard errors against Monte Carlo, the analytic Jacobian, exact population fits, the rank deficiency of a two-item specific factor, the list-experiment estimators and allocations, the mixed model and the poststratification identity, the two-step EAP shrinkage, the index quadrature, and the identities of \cref{thm:dif-bias}(d)(i) and \cref{rem:id-anchor-S}.

\paragraph{Seeds}
The base seed is $20261103$. Every replication has its own NumPy generator, created by \texttt{np.random.default\_rng} from a seed sequence \texttt{SeedSequence([20261103,\,$i$,\,keys\dots,\,$r$])}, where $i\in\{1,\dots,5\}$ is the experiment index and $r$ the replication.
\begin{itemize}
  \item E1 and E1b: keys (design index, $n$), with design indices $0,\dots,4$ for $|\mathcal I|=6$, $3$, $2$, proportional and near-proportional. The sequence is split in two, one stream for the data and one for the bootstrap and the random starts. The weights of the population fits come from one sample of $100{,}000$ respondents per design.
  \item E2: key $100\kappa$, rounded; the weights of the population fits come from samples of $60{,}000$ per group with keys $(999,\ 100\kappa)$.
  \item E3: keys (cell index), with cells numbered in the order of \cref{tab:sim-list} (parts A and B) and then \cref{tab:sim-listchoice} (part C), 32 cells in all.
  \item E4: keys $(n)$; the sequence is split into a population stream and a sample stream.
  \item E5: keys $(100s,\ 100|\Delta\mu_G|)$ for the random uniform DIF, $(700+100s_a,\ 100\Delta\mu_G)$ for the non-uniform DIF, and $(500+100s,\ 100\Delta\mu_G,\ r)$ for the finite-sample check, each rounded to an integer.
\end{itemize}
Because each replication has its own stream, the results do not depend on the number of worker processes or on the order in which tasks run.

\paragraph{Runs}
All runs were on 3~November 2026, with two worker processes, \texttt{OPENBLAS\_NUM\_THREADS=1} and \texttt{nice}, on a shared four-core machine with a load average of 11 to 14. Run times are therefore conservative.
\begin{enumerate}[(1)]
  \item E1 ran from 21:49:41 UTC, with a wall time of 1{,}276\,s and 2{,}547\,s of CPU time for 870 fits with bootstrap.
  \item The replications of E1b were computed in a run started at 22:11:23 UTC (wall time 101\,s), which later failed while summarising E2. The final run, started at 22:15:16 UTC, reused these replications unchanged (option \texttt{--reuse-raw E1b}) and recomputed E2, E3, E4 and E5 (wall times 58, 16, 5 and 22\,s).
  \item At 22:17:13 UTC the summaries of E1, E1b and E2 were recomputed from the raw replications with the final version of \texttt{run\_experiments.py} (option \texttt{--resummarize}). The E1 replications were produced by an earlier version of that script (SHA-256 prefix \texttt{ca450e64a827b7ce}). The modules E1 uses for generation and estimation (\texttt{generate.py}, \texttt{designs.py}, \texttt{polychoric.py}, \texttt{bvn.py}, \texttt{fit.py}, \texttt{boundary.py}) are byte-identical in the two versions, and only the summary code changed. E1b, E2, E3, E4 and E5 were produced by version \texttt{cd237869adbe67af}, and the summaries by the final version \texttt{2729078df210b26b}.
\end{enumerate}
The results file records, for each experiment, the SHA-256 prefixes of every source file that produced it and of the code that recomputed its summary. An earlier full run of E2 to E5 (21:41:53 UTC, in \texttt{results/run\_log.txt}) was superseded by the final run and is not used.

\subsection{Configurations}\label{app:sim-config}

\Cref{tab:sim-design} summarises the five experiments.

\begin{table}[tbp]
\centering
\caption{Design of the simulation study (synthetic data only). $R$ is the number of Monte Carlo replications.}
\label{tab:sim-design}
\footnotesize
\renewcommand{\arraystretch}{1.12}
\begin{tabularx}{\textwidth}{@{}l>{\raggedright\arraybackslash}p{3.0cm}>{\raggedright\arraybackslash}X>{\raggedright\arraybackslash}p{2.0cm}@{}}
\toprule
Exp. & Results tested & Truth and estimators & $R$ \\
\midrule
E1 & \Cref{thm:ident,prop:model-sl,prop:id-boundary-test}; \cref{rem:id-heterogeneity} &
  One group; 8 stereotype items and $|\mathcal I|\in\{2,3,6\}$ Israel-related items; heterogeneous, nearly proportional or proportional stereotype loadings; $n\in\{500,1000,3000\}$; bifactor and correlated-factors fits &
  100 ($|\mathcal I|=6,3$); 50 ($|\mathcal I|=2$); 40 (others) \\
E2 & \Cref{thm:dif-bias,prop:anchor}; \cref{rem:id-anchor-S,rem:id-concurrent} &
  Two groups of 1{,}500; uniform DIF on I1, I2 and non-uniform DIF on I3 of magnitude $\kappa\in\{0,0.2,0.4,0.6\}$; ignoring DIF, correct and contaminated anchors, stereotype-only anchors, and the DIF procedure &
  100 per $\kappa$ \\
E3 & \Cref{prop:list-var,thm:joint-sdb,thm:design} &
  One sensitive item with a threshold shift (\cref{def:le-shift}); $n=3{,}000$; split and paired designs; equal, optimal and pilot-based allocations; control lists with $J\in\{2,\dots,6\}$, with and without floor and ceiling effects &
  2{,}000 (allocation, pairing); 1{,}000 (list choice) \\
E4 & \Cref{eq:mrp,prop:eap} &
  50 areas, 32 demographic cells each, differential non-response; $n\in\{1000,3000,6000\}$; MRP, MRP on two-step EAP scores, and direct estimates &
  200 per $n$ \\
E5 & \Cref{prop:index} &
  Index $P_g(h)=\Pr(S_i\ge h\mid g)$ for the 8 stereotype items, $h\in\{2,4,6\}$; random uniform and non-uniform DIF; unequal latent variances; exact computation &
  1{,}000 configurations per cell \\
\bottomrule
\end{tabularx}
\end{table}

\paragraph{E1}
All items have five categories with standardised thresholds giving marginal shares $0.30$, $0.30$, $0.20$, $0.12$ and $0.08$, shifted by equally spaced amounts from $-0.3$ to $0.3$ across the items of a design. The link is the normal ogive. The stereotype items S1--S8 have ECVs $0.85$, $0.45$, $0.75$, $0.55$, $0.90$, $0.50$, $0.80$, $0.60$ and communalities $0.55$, $0.50$, $0.60$, $0.50$, $0.55$, $0.60$, $0.50$, $0.55$. In the proportional design every stereotype item has ECV $0.70$ and communality $0.55$. In the near-proportional design (E1b) the ECVs are $0.74$, $0.66$, $0.72$, $0.68$, $0.75$, $0.65$, $0.73$, $0.67$ with communality $0.55$. The Israel-related items I1--I6 have ECVs $0.25$, $0.75$, $0.40$, $0.60$, $0.15$, $0.85$ and communality $0.5$. Standardised loadings are $\lambda_{jG}=(\mathrm{ECV}_jh_j^2)^{1/2}$ and $\lambda_{js}=\{(1-\mathrm{ECV}_j)h_j^2\}^{1/2}$, converted to slopes and thresholds by \cref{lem:id-reduction}. Sample sizes are $n\in\{500,1000,3000\}$; $R=100$ for $|\mathcal I|=6$ and $3$, $50$ for $|\mathcal I|=2$, and $40$ for the proportional and near-proportional designs; $\tau=0.5$. The pseudo-true $\rho_{SI}$ is $0.621$, $0.589$, $0.591$, $0.590$ and $0.590$ for $|\mathcal I|=6$, $3$, $2$, proportional and near-proportional.

\paragraph{E2}
The six-item design of E1, two groups of $1{,}500$, focal latent means $(0.3,0,0.3)$ and focal covariance $\bm I_3$. Uniform DIF on I1 and I2 sets $b^{(1)}_{jc}=b_{jc}-\kappa\,(1+a_{jG}^2+a_{js}^2)^{1/2}$; non-uniform DIF on I3 sets $a^{(1)}_{jG}=(1-1.25\kappa)a_{jG}$. $\kappa\in\{0,0.2,0.4,0.6\}$, $R=100$ per $\kappa$, and the leave-one-anchor-out refits are run at $\kappa=0.4$ in all 100 replications. Anchor sets: correct S1--S4, I5, I6; one contaminated S1--S4, I1, I6; two contaminated S1--S4, I1, I2; stereotype-only S1--S4. DIF screen: Wald test of equal slopes and thresholds (6 degrees of freedom), Benjamini--Hochberg at $q=0.05$ over the tested items, effect-size threshold $0.25$.

\paragraph{E3}
The sensitive item's latent response is $0.75\,\theta_G+\varepsilon$ with $\varepsilon\sim\mathrm N(0,1)$. The private threshold gives $\pi=0.30$ and the direct threshold, shifted by $\delta_j$, gives $\pi^D=0.20$. False endorsement, when present, sets $D_i=1$ with probability $0.02/0.7$ when $Z_i=0$, so $\pi^F=0.02$. Control items have prevalence $1/2$. They form $\lfloor J/2\rfloor$ independent pairs, each with $\Pr(1,1)=\Pr(0,0)=(1+\rho)/4$, plus one independent item if $J$ is odd. $n=3{,}000$ and costs are equal. The pilot-based allocation uses a pilot of 300 respondents with equal arms, which is not reused. The covariate-adjusted estimator uses a score of $\theta_G$ with reliability $0.85$ and separate-slope regression adjustment of each arm mean, with HC2 variances. Floor and ceiling effects: in arm $\armT$, a respondent with $Z_i=1$ and $W_i=J$ reports $J$ with probability $\kappa_{\mathrm c}$, and one with $Z_i=1$ and $W_i=0$ reports $0$ with probability $\kappa_{\mathrm f}$. The theoretical MSE-optimal list is computed on the menu $J\in\{2,\dots,6\}\times\rho\in\{0,-0.6\}$ for $n\in\{300,10^3,3\times10^3,\dots,10^6\}$. $R=2{,}000$ in parts A and B and $1{,}000$ in part C.

\paragraph{E4}
Population of $10^7$ in 50 areas with sizes proportional to $\exp(Z_a)$, $Z_a\sim\mathrm N(0,1)$. Within each area, the 32 cells (age 4 $\times$ sex 2 $\times$ education 4) have a Dirichlet mix with concentration $30$ per cell around a uniform national mix. Cell means of $G$ are $b_{\mathrm{age}}+b_{\mathrm{sex}}+b_{\mathrm{edu}}+0.25\,z_a+u_a$ with $b_{\mathrm{age}}=(0.25,0.05,-0.10,-0.20)$, $b_{\mathrm{sex}}=(0.05,-0.05)$, $b_{\mathrm{edu}}=(0.20,0.10,-0.10,-0.20)$, $z_a\sim\mathrm N(0,1)$ and $u_a\sim\mathrm N(0,0.2^2)$. The within-cell standard deviation of $G$ is $0.85$. The relative response propensity is the product of an age factor $(0.6,0.8,1.1,1.4)$ and an education factor $(0.6,0.9,1.1,1.4)$. Respondents are drawn with probability proportional to cell size times propensity, and each contributes $y_i=\theta_{iG}+e_i$ with $e_i\sim\mathrm N(0,0.4^2)$. $n\in\{1000,3000,6000\}$, $R=200$, with a new population in each replication.

\paragraph{E5}
The eight stereotype items of E1. Agreement is a response in one of the top two categories. The population count distribution is computed by tensor Gauss--Hermite quadrature with 40 nodes per dimension over $(\theta_G,\theta_S)$, and the Poisson-binomial distribution of the count given the latent vector. Uniform DIF adds $\mathrm N(0,s^2)$ times the latent-response standard deviation of each item to its agreement threshold, with $s\in\{0,0.1,\dots,0.5\}$, $|\Delta\mu_G|\in\{0.05,0.1,0.2,0.3\}$ with random sign, $h\in\{2,4,6\}$ and 1{,}000 configurations per cell. Unequal variances: $\sigma_G^{(1)}\in\{0.6,0.8,1.25,1.5\}$ with $\Delta\mu_G\in\{0.1,0.2\}$. Non-uniform DIF: general slopes multiplied by $\exp\{\mathrm N(0,s_a^2)\}$, $s_a\in\{0.2,0.4,0.6\}$, $\Delta\mu_G=0.1$, 1{,}000 configurations each. Finite-sample check: $1{,}000$ respondents per group, $h=4$, $\Delta\mu_G=0.1$, $s\in\{0,0.3\}$, 300 replications each.

\subsection{Estimator settings}\label{app:sim-est}

\paragraph{Polychorics and weights}
Thresholds are normal quantiles of the cumulative marginal shares. Each polychoric correlation maximises the bivariate ordinal-probit likelihood of its table with the thresholds fixed, by vectorised Fisher scoring \citep{olsson1979}. The asymptotic covariance $\bm\Gamma$ of thresholds and correlations is the empirical second moment of per-respondent influence functions, including the effect of the estimated thresholds. The weights are $\operatorname{diag}(\bm\Gamma)^{-1}$, and the sandwich covariance is $(\bm\Delta^\top\bm W\bm\Delta)^{-1}\bm\Delta^\top\bm W\bm\Gamma\bm W\bm\Delta(\bm\Delta^\top\bm W\bm\Delta)^{-1}$, where $\bm\Delta$ is the Jacobian of the model-implied statistics. In one group only the correlations are fitted, since the thresholds are saturated. In two groups thresholds and correlations are fitted jointly, with the reference group's means set to zero and its variances to one.

\paragraph{Optimisation}
A projected Levenberg--Marquardt algorithm with analytic Jacobian, at most 200 iterations, and bounds $|a_{jf}|\le12$ (a standardised loading of $0.997$), focal log-variances in $[-3,3]$ and thresholds in $[-40,40]$. Factors are reflected so that the reference group's loadings on each factor sum to a positive number. Starts: a default ($\lambda_{jG}=0.6$, $\lambda_{js}=0.4$ for every item) and four random starts with $\lambda_{jG}\sim\mathrm U(0.2,0.8)$ and $\lambda_{js}\sim\mathrm U(0.1,0.6)$. The fit with the smallest objective is kept. The \emph{start range} of $\mathrm{ECV}_j$ is its range over the starts whose objective is within a relative $10^{-6}$ of the best. A fit is counted as \emph{at a bound} if any slope has $|a_{jf}|\ge0.999\times12$. The identification diagnostic is $\sigma_{\min}/\sigma_{\max}$ of $\bm W^{1/2}\bm\Delta$ at the estimate.

\paragraph{Intervals}
Wald intervals for $\mathrm{ECV}_j$ are formed on the logit scale from the delta-method standard error, with $z=1.96$, or $z=\Phi^{-1}\{1-0.05/(2|\mathcal I|)\}$ for the Bonferroni intervals. The bootstrap draws $B=100$ samples of respondents with replacement, re-estimates the polychorics (at most 30 scoring iterations, started at the full-sample values), and refits the bifactor model from the full-sample estimate with the full-sample weights (at most 60 iterations); percentile intervals are reported. The classification is ``in'' if the lower limit is at least $\tau$, ``out'' if the upper limit is below $\tau$, and ``undetermined'' otherwise. The correlated-factors model is fitted by the same weighted least squares with $\rho_{SI}=\tanh z$.

\paragraph{E3 and E4}
The list estimators are those of \cref{thm:joint-sdb}, with plug-in variances from \eqref{eq:le-split}, \eqref{eq:le-vps} and \eqref{eq:le-vh}. The combined estimator is set to $0$ if a stratum $D_i=0$ of a list arm has fewer than two respondents. MRP uses restricted maximum likelihood for the ratio of the area variance to the residual variance (bounded one-dimensional search on its logarithm), generalised least squares for the fixed effects, best linear unbiased predictors for the area effects, and the prediction-error variance from Henderson's mixed-model equations. The two-step scores use the prior $\mathrm N(m_0,\sigma_0^2)$ with $m_0=\bar y$ and $\sigma_0^2=\widehat{\Var}(y)-0.4^2$, so that $\hat\varrho=\sigma_0^2/(\sigma_0^2+0.4^2)$.

\subsection{Monte Carlo standard errors}\label{app:sim-mcse}

For a mean the MCSE is the replication SD divided by $\sqrt R$; for a proportion $p$ it is $\{p(1-p)/R\}^{1/2}$; for an RMSE it is the delta-method value $\mathrm{SD}(e^2)/(2\,\mathrm{RMSE}\sqrt R)$; and for a median it is $1.2533\times\mathrm{IQR}/(1.349\sqrt R)$. A coverage near 95\% therefore has an MCSE of about $2.2$ percentage points with $R=100$ and $3.4$ with $R=40$. Pooled E1 quantities treat the $|\mathcal I|$ items of a replication as independent, which understates their MCSE. In E2 the MCSE of the bias of $\hat\mu_G$ is $0.005$ to $0.020$. In E3 the MCSE of the bias is $0.0005$ to $0.0017$. In E4 the MCSE of the RMSE is at most $0.004$. In E5 a reversal rate from 1{,}000 configurations has an MCSE of at most $1.6$ percentage points. All MCSEs are stored in the results file next to the quantities they refer to.

\subsection{Exact checks and recomputed illustrations}\label{app:sim-checks}

Three identities of \cref{sec:identification} were checked on exact population statistics in E2. Uniform DIF $-\bm a_j^\top\bm v$ on every item, with $\bm v=(0.25,-0.10,0.15)$, is fitted exactly by the invariant model with latent mean $\bm\mu+\bm v=(0.55,-0.10,0.45)$ (\cref{thm:dif-bias}(d)(i)). Three anchors (S1, S2, I1) with DIF $-\bm A_{\mathcal A}\bm v$ give the same shift, and every non-anchor item acquires spurious DIF equal to $\bm a_k^\top\bm v$ to within $2.3\times10^{-9}$ (\cref{prop:anchor}(b)). Stereotype-only anchors recover $\mu_G$ and $\mu_S$ to within $2\times10^{-13}$ and every focal ECV to within $2.1\times10^{-14}$ under the Israel-related DIF of E2 at $\kappa=0.4$ (\cref{rem:id-anchor-S}).

The results file also recomputes, independently, the illustrative numbers quoted in \cref{sec:sensitive,sec:smallarea,sec:design}. For four independent control items with prevalence $1/2$ and $\pi=0.2$, $L^*=26.96$; for two pairs with $\rho=-0.6$, $L^*=11.92$ (\cref{prop:list-var}). The thresholds of \cref{rem:le-mse} are $0.0456$ ($n=2{,}000$) and $0.0263$ ($n=6{,}000$). For \cref{thm:design}(a) with $\sigma_{\armC}^2=\sigma_{\armT}^2=1$ and $\pi^D=0.2$, $n\Var$ is $6.48$ under equal allocation and $5.76$ at the optimum, with shares $(1/6,5/12,5/12)$. The direct-arm cost threshold of \cref{thm:design}(b) is $1.26\times10^{-4}$, and the paired-to-split variance ratio of \cref{cor:le-paired-split} is $0.810$. The arm size for the power calculation of \cref{sec:design-list} is $1{,}695$ for $\SDB=0.10$ and $6{,}781$ for $\SDB=0.05$. \Cref{ex:sa-index} gives $P_1(5)=0.100$ and $P_2(5)=0.214$ in (i), $P_1(1)=0.999$, $P_2(1)=0.806$, $P_1(10)=0.002$ and $P_2(10)=0.194$ in (ii), and changes of $0.0049$ and $0.0987$ in (iii).

\subsection{Further result tables}\label{app:sim-tables}

\Cref{tab:sim-list,tab:sim-listchoice,tab:sim-mrp,tab:sim-index} give the full results of E3, E4 and E5 that \cref{sec:sim-list,sec:sim-mrp,sec:sim-index} summarise.

\begin{table}[tbp]
\centering
\caption{E3 parts A and B, estimation of social-desirability bias (synthetic data; $n=3{,}000$; $R=2{,}000$ per row). Shares $(r_D,r_C,r_T)$ of the direct, list-control and list-treatment arms (means over replications for the pilot-based allocation). SD (theory): exact standard deviation from \eqref{eq:le-split}, or asymptotic from \eqref{eq:le-vps} and \eqref{eq:le-vh}, at the allocation used. Bias and RMSE are as estimators of $\SDB$, which is $0.10$ when $\pi^F=0$ and $0.08$ when $\pi^F=0.02$. Cov.: coverage (\%) of the 95\% Wald interval with the plug-in standard error; SE/SD: mean standard error over the Monte Carlo SD.}
\label{tab:sim-list}
\scriptsize
\setlength{\tabcolsep}{4pt}
% Generated by code/make_figures.py from results/sim_results.json -- do not edit by hand.
% X3 parts A, B: n = 3000, pi = 0.3, SDB = 0.10 (pi^F = 0) or SDB = 0.10 - pi^F; R = 2000 per row. Bias/RMSE of each estimator as an estimator of SDB.
% SD (theory): exact (eq:le-split) or asymptotic (eq:le-vps, eq:le-vh) SD at the allocation used;
% Cov.: coverage (%) of the 95% Wald interval with the plug-in SE; SE/SD: mean SE over Monte Carlo SD.
\begin{tabular}{@{}llrrrrrrr@{}}
\toprule
Design and allocation & Estimator & $(r_D,r_C,r_T)$ & SD (theory) & Bias & SD & RMSE & Cov. & SE/SD \\
\midrule
\multicolumn{9}{@{}l}{\emph{A. Split design, $J=4$, $\pi^F=0$}} \\
indep.\ controls, equal & $\widehat{\mathrm{SDB}}_{\mathrm{S}}$ & (0.33, 0.33, 0.33) & 0.0487 & +0.0010 & 0.0484 & 0.0484 & 95.5 & 1.00 \\
indep.\ controls, Neyman & $\widehat{\mathrm{SDB}}_{\mathrm{S}}$ & (0.16, 0.40, 0.44) & 0.0456 & +0.0008 & 0.0459 & 0.0459 & 94.3 & 0.99 \\
indep.\ controls, pilot plug-in & $\widehat{\mathrm{SDB}}_{\mathrm{S}}$ & (0.16, 0.40, 0.44) & 0.0457 & +0.0016 & 0.0462 & 0.0462 & 95.1 & 0.99 \\
pairs $\rho=-0.6$, equal & $\widehat{\mathrm{SDB}}_{\mathrm{S}}$ & (0.33, 0.33, 0.33) & 0.0342 & $-$0.0007 & 0.0341 & 0.0341 & 95.2 & 1.00 \\
pairs $\rho=-0.6$, Neyman & $\widehat{\mathrm{SDB}}_{\mathrm{S}}$ & (0.22, 0.35, 0.43) & 0.0331 & +0.0008 & 0.0332 & 0.0332 & 95.2 & 1.00 \\
pairs $\rho=-0.6$, pilot plug-in & $\widehat{\mathrm{SDB}}_{\mathrm{S}}$ & (0.22, 0.35, 0.43) & 0.0332 & +0.0008 & 0.0324 & 0.0324 & 95.2 & 1.02 \\
\addlinespace
\multicolumn{9}{@{}l}{\emph{B. Paired vs split, $J=4$, pairs $\rho=-0.6$}} \\
split, Neyman, $\pi^F=0$ & $\widehat{\mathrm{SDB}}_{\mathrm{S}}$ & (0.22, 0.35, 0.43) & 0.0331 & +0.0002 & 0.0339 & 0.0339 & 94.4 & 0.98 \\
paired, optimal, $\pi^F=0$ & $\widehat{\mathrm{SDB}}_{\mathrm{PS}}$ & (0.00, 0.47, 0.53) & 0.0243 & $-$0.0002 & 0.0249 & 0.0249 & 94.5 & 0.98 \\
 & $\widehat\pi^H$ & (0.00, 0.47, 0.53) & 0.0220 & +0.0000 & 0.0226 & 0.0226 & 94.8 & 0.97 \\
paired, equal, $\pi^F=0$ & $\widehat{\mathrm{SDB}}_{\mathrm{PS}}$ & (0.33, 0.33, 0.33) & 0.0298 & $-$0.0002 & 0.0296 & 0.0296 & 95.5 & 1.01 \\
 & $\widehat\pi^H$ & (0.33, 0.33, 0.33) & 0.0270 & $-$0.0001 & 0.0268 & 0.0268 & 95.5 & 1.01 \\
split, Neyman, $\pi^F=0.02$ & $\widehat{\mathrm{SDB}}_{\mathrm{S}}$ & (0.23, 0.35, 0.43) & 0.0334 & +0.0013 & 0.0337 & 0.0338 & 94.5 & 0.99 \\
paired, optimal, $\pi^F=0.02$ & $\widehat{\mathrm{SDB}}_{\mathrm{PS}}$ & (0.00, 0.47, 0.53) & 0.0246 & +0.0007 & 0.0242 & 0.0242 & 95.2 & 1.02 \\
 & $\widehat\pi^H$ & (0.00, 0.47, 0.53) & 0.0218 & +0.0206 & 0.0214 & 0.0297 & 85.3 & 1.01 \\
paired, equal, $\pi^F=0.02$ & $\widehat{\mathrm{SDB}}_{\mathrm{PS}}$ & (0.33, 0.33, 0.33) & 0.0301 & $-$0.0002 & 0.0302 & 0.0302 & 95.0 & 1.00 \\
 & $\widehat\pi^H$ & (0.33, 0.33, 0.33) & 0.0267 & +0.0198 & 0.0263 & 0.0329 & 89.8 & 1.01 \\
\bottomrule
\end{tabular}

\end{table}

\begin{table}[tbp]
\centering
\caption{E3 part C, choice of control list (synthetic data; split design, Neyman allocation, $n=3{,}000$; $R=1{,}000$ per row). Exposure is $\Pr(W_i\in\{0,J\})$. Theoretical bias from \eqref{eq:le-fcbias} and SD from \eqref{eq:le-split}; Cov.: coverage (\%) of the 95\% Wald interval.}
\label{tab:sim-listchoice}
\footnotesize
% Generated by code/make_figures.py from results/sim_results.json -- do not edit by hand.
% E3 part C: split design, Neyman allocation, n = 3000, R = 1000 per row; kappa = kc = kf (floor/ceiling); exposure = Pr(W in {0, J}).
% MSE-optimal list by n (thm:design(c), theory): kappa=0: {'300': 'J2_rho-0.6', '1000': 'J2_rho-0.6', '3000': 'J2_rho-0.6', '10000': 'J2_rho-0.6', '30000': 'J2_rho-0.6', '100000': 'J2_rho-0.6', '300000': 'J2_rho-0.6', '1000000': 'J2_rho-0.6'}; kappa=0.5: {'300': 'J2_rho-0.6', '1000': 'J2_rho-0.6', '3000': 'J4_rho-0.6', '10000': 'J4_rho-0.6', '30000': 'J4_rho-0.6', '100000': 'J4_rho-0.6', '300000': 'J6_rho-0.6', '1000000': 'J6_rho-0.6'}.
% Empirical MSE-optimal list at n = 3000: kappa=0: J2_rho-0.6; kappa=0.5: J4_rho-0.6.
\begin{tabular}{@{}rrrrrrrrrr@{}}
\toprule
 & & & & \multicolumn{2}{c}{Bias} & \multicolumn{2}{c}{SD} & & \\
\cmidrule(lr){5-6} \cmidrule(lr){7-8}
$J$ & $\rho$ & $\kappa_{\mathrm{c}}=\kappa_{\mathrm{f}}$ & Exposure & Theory & Sim. & Theory & Sim. & RMSE & Cov. \\
\midrule
2 & 0 & 0 & 0.500 & +0.0000 & +0.0000 & 0.0356 & 0.0367 & 0.0367 & 94.6 \\
2 & $-$0.6 & 0 & 0.200 & +0.0000 & $-$0.0005 & 0.0272 & 0.0283 & 0.0283 & 94.0 \\
3 & 0 & 0 & 0.250 & +0.0000 & $-$0.0007 & 0.0410 & 0.0421 & 0.0421 & 94.2 \\
3 & $-$0.6 & 0 & 0.100 & +0.0000 & +0.0008 & 0.0344 & 0.0347 & 0.0347 & 95.0 \\
4 & 0 & 0 & 0.125 & +0.0000 & +0.0004 & 0.0456 & 0.0468 & 0.0468 & 94.2 \\
4 & $-$0.6 & 0 & 0.020 & +0.0000 & $-$0.0008 & 0.0331 & 0.0331 & 0.0331 & 96.1 \\
5 & 0 & 0 & 0.062 & +0.0000 & +0.0007 & 0.0498 & 0.0500 & 0.0500 & 95.6 \\
5 & $-$0.6 & 0 & 0.010 & +0.0000 & +0.0027 & 0.0390 & 0.0391 & 0.0392 & 95.0 \\
6 & 0 & 0 & 0.031 & +0.0000 & +0.0001 & 0.0535 & 0.0520 & 0.0520 & 95.3 \\
6 & $-$0.6 & 0 & 0.002 & +0.0000 & $-$0.0001 & 0.0379 & 0.0355 & 0.0354 & 96.4 \\
\addlinespace
2 & 0 & 0.5 & 0.500 & $-$0.0750 & $-$0.0726 & 0.0356 & 0.0361 & 0.0811 & 44.1 \\
2 & $-$0.6 & 0.5 & 0.200 & $-$0.0300 & $-$0.0313 & 0.0272 & 0.0272 & 0.0415 & 78.4 \\
3 & 0 & 0.5 & 0.250 & $-$0.0375 & $-$0.0365 & 0.0410 & 0.0419 & 0.0556 & 85.8 \\
3 & $-$0.6 & 0.5 & 0.100 & $-$0.0150 & $-$0.0145 & 0.0344 & 0.0348 & 0.0377 & 92.4 \\
4 & 0 & 0.5 & 0.125 & $-$0.0187 & $-$0.0180 & 0.0456 & 0.0440 & 0.0475 & 95.0 \\
4 & $-$0.6 & 0.5 & 0.020 & $-$0.0030 & $-$0.0033 & 0.0331 & 0.0329 & 0.0331 & 94.4 \\
5 & 0 & 0.5 & 0.062 & $-$0.0094 & $-$0.0084 & 0.0498 & 0.0503 & 0.0510 & 94.7 \\
5 & $-$0.6 & 0.5 & 0.010 & $-$0.0015 & $-$0.0007 & 0.0390 & 0.0384 & 0.0383 & 95.2 \\
6 & 0 & 0.5 & 0.031 & $-$0.0047 & $-$0.0045 & 0.0535 & 0.0528 & 0.0530 & 95.3 \\
6 & $-$0.6 & 0.5 & 0.002 & $-$0.0003 & +0.0013 & 0.0379 & 0.0371 & 0.0371 & 95.2 \\
\bottomrule
\end{tabular}

\end{table}

\begin{table}[tbp]
\centering
\caption{E4, area means of $G$ (synthetic data; 50 areas; $R=200$ populations and samples per $n$). Empty: mean number of areas without respondents. Cov.: coverage (\%) of 95\% intervals. Slope: regression of estimated on true area deviations from the national mean (1 means no attenuation). Direct estimators are evaluated on areas with at least two respondents.}
\label{tab:sim-mrp}
\footnotesize
\setlength{\tabcolsep}{4.5pt}
% Generated by code/make_figures.py from results/sim_results.json -- do not edit by hand.
% X4: 50 areas; R = 200 populations x surveys per n. Direct estimators evaluated on areas with >= 2 respondents; MRP on all areas.
% Cov.: coverage (%) of 95% intervals (Henderson prediction-error variance, variance components plugged in).
% Slope: regression of estimated on true area deviations from the national mean (1 = no attenuation);
% MRP's slope < 1 is partial pooling; the two-step EAP slope is smaller by the factor rho of prop:eap.
\begin{tabular}{@{}rrrrrrrrrrr@{}}
\toprule
 & & \multicolumn{3}{c}{MRP} & \multicolumn{3}{c}{MRP on EAP scores} & \multicolumn{2}{c}{Direct, unweighted} & Direct, wtd. \\
\cmidrule(lr){3-5} \cmidrule(lr){6-8} \cmidrule(lr){9-10} \cmidrule(lr){11-11}
$n$ & Empty & RMSE & Cov. & Slope & RMSE & Cov. & Slope & RMSE & Bias & RMSE \\
\midrule
1000 & 0.8 & 0.168 & 92.0 & 0.75 & 0.172 & 85.3 & 0.64 & 0.342 & $-$0.099 & 0.359 \\
3000 & 0.1 & 0.133 & 93.7 & 0.83 & 0.141 & 86.0 & 0.70 & 0.232 & $-$0.098 & 0.229 \\
6000 & 0.0 & 0.105 & 94.7 & 0.90 & 0.115 & 84.0 & 0.76 & 0.177 & $-$0.094 & 0.164 \\
\bottomrule
\end{tabular}

\end{table}

\begin{table}[tbp]
\centering
\caption{E5, sign reversals of stereotype-count indices (synthetic data; exact population computations). Percentage of 1{,}000 random uniform-DIF configurations per cell in which $P_1(h)-P_0(h)$ has the opposite sign to $\Delta\mu_G$. Columns: $|\Delta\mu_G|$ within each index threshold $h$ (of $M=8$); rows: standard deviation $s$ of the threshold shifts, in standard deviations of each item's latent response.}
\label{tab:sim-index}
\footnotesize
\setlength{\tabcolsep}{4.5pt}
% Generated by code/make_figures.py from results/sim_results.json -- do not edit by hand.
% E5: percentage of random uniform-DIF configurations (1000 per cell) in which the population index difference P1(h) - P0(h) has the opposite sign to Delta mu_G = mu_G^(1) - mu_G^(0);
% columns: |Delta mu_G| within each h; rows: SD s of the DIF shifts (in SDs of the latent response).
% Random non-uniform DIF (G slopes x exp(N(0, s_a^2))), Delta mu_G = 0.1: s_a=0.2: h-dependent sign 25.9%; s_a=0.4: h-dependent sign 48.3%; s_a=0.6: h-dependent sign 63.6%.
\begin{tabular}{@{}rrrrrrrrrrrrr@{}}
\toprule
 & \multicolumn{4}{c}{$h=2$} & \multicolumn{4}{c}{$h=4$} & \multicolumn{4}{c}{$h=6$} \\
\cmidrule(lr){2-5} \cmidrule(lr){6-9} \cmidrule(lr){10-13}
$s$ & $0.05$ & $0.1$ & $0.2$ & $0.3$ & $0.05$ & $0.1$ & $0.2$ & $0.3$ & $0.05$ & $0.1$ & $0.2$ & $0.3$ \\
\midrule
0 & 0.0 & 0.0 & 0.0 & 0.0 & 0.0 & 0.0 & 0.0 & 0.0 & 0.0 & 0.0 & 0.0 & 0.0 \\
0.1 & 21.1 & 5.6 & 0.0 & 0.0 & 20.1 & 5.1 & 0.0 & 0.0 & 20.1 & 5.2 & 0.0 & 0.0 \\
0.2 & 33.5 & 21.1 & 6.5 & 0.7 & 32.4 & 18.9 & 5.5 & 0.3 & 33.0 & 21.4 & 5.5 & 0.3 \\
0.3 & 41.3 & 30.2 & 14.6 & 7.8 & 39.7 & 28.4 & 12.7 & 4.6 & 39.4 & 28.7 & 14.5 & 4.2 \\
0.4 & 44.4 & 36.7 & 23.5 & 14.8 & 41.6 & 36.8 & 18.3 & 9.4 & 43.0 & 37.6 & 20.8 & 10.2 \\
0.5 & 47.4 & 41.1 & 32.1 & 24.5 & 46.2 & 38.3 & 25.2 & 17.5 & 46.6 & 39.3 & 25.8 & 18.1 \\
\bottomrule
\end{tabular}

\end{table}

\subsection{Reproduction}\label{app:sim-repro}

From \texttt{code/}, the commands
{\small
\begin{verbatim}
export OPENBLAS_NUM_THREADS=1
nice -n 10 python run_experiments.py --only E1
nice -n 10 python run_experiments.py --only E1b E2 E3 E4 E5
python run_experiments.py --resummarize E1 E1b E2
python make_figures.py
python -m pytest tests/
\end{verbatim}
}
regenerate \texttt{results/sim\_results.json}, \texttt{results/raw/}, the figures \texttt{figures/fig-sim-*.pdf}, the tables \texttt{tables/tab-sim-*.tex} and the test report. A single call without \texttt{--only} runs all experiments, and \texttt{--quick} runs a small smoke test. On the shared machine described above, the full study takes about 25 minutes of wall time with two workers, almost all of it in E1. The figures take a few seconds and the tests about six seconds. Because each replication has its own seed, the numbers do not depend on the number of workers. Results can differ in the last digits across platforms or library versions, through floating-point differences in the optimiser.
